\section{Experimental Setup}
\label{sec:experiments}

This section specifies the data, the evaluation protocol, and the
configurations compared. The numerical results are deferred to
\Cref{sec:results}.

\subsection{Datasets}

\subsubsection*{Training data}
Training combines two sources, merged into a single instruction-tuned JSONL:
\begin{itemize}
    \item \textbf{AnnoMI alcohol-only subset.} Filtered from the full AnnoMI
    corpus~\cite{wu2022annomi} to dialogues whose declared topic is alcohol use,
    converted into instruction/response pairs for SFT. After preprocessing,
    363 training examples and 40 held-out evaluation examples remain.
    \item \textbf{Synthetic MI conversations.} Generated by the pipeline
    described in \Cref{sec:analysis} and stored as 13 batch JSONL files under
    \texttt{data/synthetic/}, totalling approximately one thousand
    conversations across short, medium, long, and mixed-defensiveness
    scenario types.
\end{itemize}

\subsubsection*{Knowledge base for retrieval}
The retrieval corpus is described in \Cref{tab:kb-contents}. It is held strictly
separate from the training corpus, so retrieved content is not memorized text
the model has already seen during fine-tuning.

\subsubsection*{Evaluation set}
The evaluation set contains 181 user turns drawn from three sources:
\begin{itemize}
    \item $\sim$80 real AnnoMI turns (held out from training),
    \item $\sim$80 synthetic eval turns spanning emotion $\times$ defensiveness
    combinations,
    \item $\sim$20 hand-crafted edge cases (deeply ambivalent, sarcastic,
    minimizing, and crisis-adjacent phrasings).
\end{itemize}
Each item provides a user turn, a reference response (where available), and
metadata for analysis (emotion, defensiveness, talk type).

\subsection{Configurations Compared}

Five configurations are evaluated on the same set of user turns. Each
configuration produces a single response per turn, scored independently by
the judge.

\begin{table}[h]
\centering
\caption{Configurations under evaluation. RAG denotes retrieval over the MI
knowledge base described in \Cref{tab:kb-contents}.}
\label{tab:configs}
\begin{tabular}{lll}
\toprule
\textbf{Tag} & \textbf{Model} & \textbf{Augmentation} \\
\midrule
\texttt{base}    & \texttt{gemma-2-2b-it} (no fine-tune) & system prompt only \\
\texttt{v1}      & Gemma-2 LoRA, AnnoMI only             & system prompt only \\
\texttt{v2}      & Gemma-2 LoRA, AnnoMI + synthetic      & system prompt only \\
\texttt{v2\_rag} & Gemma-2 LoRA, AnnoMI + synthetic      & system prompt + RAG \\
\texttt{claude}  & Claude Haiku via Anthropic API        & system prompt only \\
\bottomrule
\end{tabular}
\end{table}

\texttt{base} is the no-fine-tune ablation; \texttt{v1} isolates the effect of
adding fine-tuning on real-only data; \texttt{v2} isolates the effect of
adding synthetic data; \texttt{v2\_rag} adds retrieval on top of \texttt{v2};
and \texttt{claude} is a much larger reference model that gives an upper-bound
sense of what is achievable on the same task without size constraints.

\subsection{Evaluation Protocol}

Each configuration's response is independently scored by Claude Sonnet
(\texttt{claude-sonnet-4-6}) acting as judge, using the rubric in
\Cref{tab:rubric}. The judge is given the user turn, the model response, and
the reference response (when available), and is instructed to produce a
1--5 integer score on each of five dimensions, plus a free-text rationale and
a tag for the detected MI technique.

\begin{table}[h]
\centering
\caption{Five-dimension MI rubric used by the judge model. Each dimension is
scored 1--5, where 1 = absent or counterproductive and 5 = expert-level.}
\label{tab:rubric}
\begin{tabularx}{\linewidth}{l X}
\toprule
\textbf{Dimension} & \textbf{What it measures} \\
\midrule
\textit{Empathy}              & Does the response acknowledge the user's emotional state without minimizing or pathologizing it? \\
\textit{MI technique}         & Does the response apply OARS (open question, affirmation, reflection, summary) appropriately for the turn? \\
\textit{Resistance handling}  & When the user is defensive, does the response roll with resistance instead of arguing? \\
\textit{Autonomy support}     & Does the response leave choices to the user rather than directing them? \\
\textit{Appropriateness}      & Is the response calibrated in length, tone, and content for the user's stated context? \\
\bottomrule
\end{tabularx}
\end{table}

\subsubsection*{Why Claude Sonnet as judge}
Sonnet is from a different provider than the systems under test except for
\texttt{claude} (Haiku); this minimizes the most pathological self-preference
bias~\cite{chiang2023closer}. The Haiku reference is therefore the
configuration most exposed to bias, which is acceptable because it is being
used as an upper-bound reference rather than a configuration competing for the
``best'' label.

\subsubsection*{What is reported}
For each configuration, mean scores per dimension and an overall mean
(unweighted average across dimensions) are reported. Sample standard
deviations across the 181 items are reported in \Cref{sec:results} where
relevant. Item-level technique distributions are also reported, since they
characterize \emph{what} the model is doing differently, not just
\emph{how well} it scored.
